\chapter{Intersection and Distance of Convex Objects}
\label{ch:convex_intersection}

\section{Convex Polygon Intersection}
\label{sec:convexint_polygon_intersection}

\subsection{1. Overview}

Given two convex polygons $P$ and $Q$ in the plane, the intersection problem
asks whether $P \cap Q$ is empty and, if not, returns the convex polygon
$P \cap Q$. Because the inputs are convex, the intersection is either empty
or a convex polygon with at most $|P| + |Q|$ vertices, and it can be
computed in linear time. The general (non-convex) polygon intersection
problem is harder and is handled by the segment-intersection and overlay
machinery of Chapter~\ref{ch:polygon_ops}; convexity is what makes this
chapter's results possible.

\subsection{2. Definitions}

\begin{definition}[Polygon Intersection]\label{def:convexint_polygon_inter}
  The \textbf{intersection} of two convex polygons $P$ and $Q$ is the set
  $P \cap Q$, which is empty or a convex polygon. Its vertices are either
  vertices of $P$ or $Q$ inside the other polygon, or crossing points of
  the two boundaries.
\end{definition}

\subsection{3. Theory}

Two classic facts underlie the linear-time algorithms. First, since the
boundaries of two convex polygons cross at most $|P|+|Q|$ times, the
intersection polygon has $O(|P|+|Q|)$ vertices. Second, the two boundaries
can be merged by a plane sweep in a single pass because each boundary is
sorted around its hull; this is the basis of the $O(n)$ algorithm of Shamos
\cite{preparata1985}. The distance version of the same problem --- finding
the closest points of two convex polygons --- is solved in $O(\log n)$
after linear preprocessing by the Dobkin--Kirkpatrick hierarchy
\cite{dobkin1983,dobkin1985}, and in $O(n)$ by rotating calipers
\cite{toussaint1983}.

\subsection{4. Pseudocode}

\begin{algorithm}[htbp]
\caption{Intersection of Two Convex Polygons}
\label{alg:convexint_polygon_intersection}
\begin{algorithmic}[1]
\Require Two convex polygons $P, Q$ with vertices in counterclockwise order.
\Ensure The convex polygon $P \cap Q$, or empty.
\Function{ConvexPolygonIntersection}{$P$, $Q$}
  \State $R \gets \emptyset$; $i \gets 0$; $j \gets 0$; $n \gets |P|$; $m \gets |Q|$
  \While{$i < n$ or $j < m$}
    \State test edge $e_P = (P_i, P_{i+1})$ against edge $e_Q = (Q_j, Q_{j+1})$ using orient()
    \State add crossing points of $e_P, e_Q$ to $R$
    \State advance the edge whose supporting line the other polygon is ``ahead of'' ($i$ or $j$)
    \State add the endpoint that lies inside the other polygon to $R$
  \EndWhile
  \State output the vertices of $R$ in order as the intersection polygon
\EndFunction
\end{algorithmic}
\end{algorithm}

\subsection{5. Parameter Selection}

\begin{itemize}
\item \textbf{Orientation convention}: both polygons must be traversed in
  the same orientation (counterclockwise); a clockwise input must be
  reversed first.
\item \textbf{Degeneracies}: touching edges (one vertex on an edge of the
  other) and coincident edges must be handled by the exact predicate
  discipline of Section~\ref{sec:robust_predicates}, and reported once.
\end{itemize}

\subsection{6. Complexity}

\begin{itemize}
\item \textbf{Time}: $O(n + m)$.
\item \textbf{Space}: $O(n + m)$ for the output (which has at most
  $n + m$ vertices).
\end{itemize}

\subsection{7. Usages}

\begin{itemize}
\item \textbf{Boolean geometry}: union/intersection/difference of convex
  parts (Chapter~\ref{ch:polygon_ops}).
\item \textbf{Windowing and culling}: computing the visible overlap of two
  convex windows.
\end{itemize}

\subsection{8. Testing Methods and Edge Cases}

\begin{itemize}
\item \textbf{Separation}: disjoint polygons with large and small gaps.
\item \textbf{Containment}: $Q \subset P$, $P \subset Q$, and identical
  polygons.
\item \textbf{Boundary contact}: polygons sharing an edge, a vertex, or a
  vertex on an edge.
\end{itemize}

\subsection{9. References}

The linear intersection algorithm is due to Shamos (1978 thesis) and is
presented in \cite{preparata1985}; the logarithmic distance and
intersection queries follow from the Dobkin--Kirkpatrick hierarchy
\cite{dobkin1983,dobkin1985} and the rotating-calipers technique of
Toussaint \cite{toussaint1983}.

\section{Separating Axis Theorem}
\label{sec:convexint_sat}

\subsection{1. Overview}

The \textbf{separating axis theorem} (SAT)\index{separating axis theorem}
gives a finite test for whether two convex polygons (in 2D) or convex
polyhedra (in 3D) intersect: two convex sets are disjoint exactly when
there is an axis such that their projections onto that axis are disjoint.
For polygons and polyhedra, only a small finite set of axes --- the face
normals and, in 3D, the cross products of edge directions --- must be
tested. SAT is the fastest intersection test for boxes and for convex
primitives in physics engines, and it underlies the separating plane
reasoning of Section~\ref{sec:math_separating_surfaces}.

\subsection{2. Definitions}

\begin{definition}[Projection Interval]\label{def:convexint_projection}
  The \textbf{projection} of a convex set $C$ onto an axis $u$ is the
  interval $[\min_{x\in C} \langle x, u\rangle, \max_{x\in C} \langle x,
  u\rangle]$. Two sets are \textbf{separated by axis $u$} if their
  projection intervals are disjoint.
\end{definition}

\subsection{3. Theory}

\begin{theorem}[Separating Axis Theorem]\label{thm:convexint_sat}
  Two convex polygons in the plane are disjoint iff their projections are
  disjoint on some axis perpendicular to one of their edges. Two convex
  polyhedra in $\mathbb{R}^3$ are disjoint iff their projections are
  disjoint on some axis perpendicular to one of their faces, or parallel to
  the cross product of one edge of each polyhedron.
\end{theorem}

The theorem follows from the hyperplane-separation theorem of
Section~\ref{sec:math_separating_surfaces}: the separating line (or plane)
can be rotated until it supports both sets, at which point it is parallel
to an edge (or a face, or an edge-edge cross product). For two
axis-aligned boxes, the candidate axes are the three coordinate axes, which
makes the box test three interval comparisons.

\subsection{4. Pseudocode}

\begin{algorithm}[htbp]
\caption{Separating Axis Test for Two Convex Polygons}
\label{alg:convexint_sat}
\begin{algorithmic}[1]
\Require Convex polygons $P$, $Q$.
\Ensure True iff $P \cap Q = \emptyset$.
\Function{SATIntersect}{$P$, $Q$}
  \For{each edge $e$ of $P$ \textbf{and} each edge $e$ of $Q$}
    \State $u \gets$ unit normal of $e$
    \State $I_P \gets$ projection of $P$ onto $u$; $I_Q \gets$ projection of $Q$ onto $u$
    \If{$I_P \cap I_Q = \emptyset$} \State \Return \textsc{False} \EndIf
  \EndFor
  \State \Return \textsc{True}
\EndFunction
\end{algorithmic}
\end{algorithm}

\subsection{5. Parameter Selection}

\begin{itemize}
\item \textbf{Test order}: in practice, short-circuit on the most
  discriminating axis first (e.g., the bounding-box axes), which exits early
  for most disjoint pairs.
\item \textbf{3D}: the candidate set is all face normals plus all
  edge-pair cross products; for boxes this is 6 normals (only 3 distinct)
  plus 9 cross products, but only 3 are distinct.
\end{itemize}

\subsection{6. Complexity}

\begin{itemize}
\item \textbf{Time}: $O(n + m)$ with $n, m$ the numbers of vertices; for
  boxes, $O(1)$.
\item \textbf{Space}: $O(1)$.
\end{itemize}

\subsection{7. Usages}

\begin{itemize}
\item \textbf{Physics engines}: box-box and box-convex tests in real-time
  physics use SAT because of its tiny constant and early exit.
\item \textbf{Culling}: as a fast filter before the precise (GJK) test of
  Section~\ref{sec:convexint_gjk}.
\end{itemize}

\subsection{8. Testing Methods and Edge Cases}

\begin{itemize}
\item \textbf{Touching configurations}: verify that touching (projections
  sharing an endpoint) is treated consistently as intersecting or not.
\item \textbf{Degenerate polygons}: collinear vertices produce zero-length
  edges and must be filtered before computing normals.
\end{itemize}

\subsection{9. References}

SAT is folklore in the game-physics community and is developed carefully in
\cite{ericson2004,vandenbergen2004}; its theoretical basis is the
separation theorem of Section~\ref{sec:math_separating_surfaces} and
\cite{deberg2008}.

\section{Distance between Convex Sets}
\label{sec:convexint_distance}

The \textbf{distance} between two convex sets $A$ and $B$ is
$\operatorname{dist}(A,B) = \min_{a\in A, b\in B} \|a-b\|$. The minimum is
attained because the sets are closed; the minimizers are the \textbf{closest
points} of $A$ and $B$. For convex polygons in the plane, the closest pair
is found in $O(n)$ time by rotating calipers \cite{toussaint1983}, and in
$O(\log n)$ time after $O(n)$ preprocessing by the Dobkin--Kirkpatrick
method \cite{dobkin1985}. For convex polyhedra in $\mathbb{R}^3$ the same
hierarchy gives an $O(\log^2 n)$ distance query, while the iterative
GJK algorithm of Section~\ref{sec:convexint_gjk} computes the distance by
support-function evaluations without any preprocessing.

Distance computations between convex sets are the primitive of collision
response (how far must the objects move apart?), of motion planning
(Chapter~\ref{ch:motion_planning}), and of the minimum-separation checks in
Section~\ref{sec:convexint_penetration_depth}. Two structural facts are
used throughout:

\begin{itemize}
  \item the distance between $A$ and $B$ equals the distance from the
    origin to the Minkowski difference $A \ominus B$ defined in
    Section~\ref{sec:convexint_minkowski_difference};
  \item the closest points are found by projecting onto the supporting
    plane of the nearest feature (vertex, edge, or face), which is exactly
    the geometric intuition exploited by GJK and by the Gilbert--Johnson--
    Keerthi family of algorithms.
\end{itemize}

\index{distance between convex sets}\index{closest points}

\section{Support Functions}
\label{sec:convexint_support_functions}

\begin{definition}[Support Function and Support Point]\label{def:convexint_support}
  For a convex set $C$ and a direction $u$, the \textbf{support point}
  $s_C(u)$ is a point of $C$ maximizing $\langle x, u\rangle$ over $x \in
  C$. The \textbf{support function} $h_C(u) = \langle s_C(u), u\rangle$ is
  the maximum of the linear functional in direction $u$.
\end{definition}

Support points are the fundamental oracle for convex intersection and
distance algorithms, because many questions about $C$ reduce to evaluating
$s_C$ in a few directions and $s_C$ is easy to compute for every convex
primitive:
\begin{itemize}
  \item \textbf{box}: $s(u) = (\operatorname{sign}(u_x)L_x, \dots)$, one
    comparison per coordinate;
  \item \textbf{polygon}: $O(\log n)$ by binary search over the edge
    directions, or $O(n)$ by scanning;
  \item \textbf{sphere/ellipsoid}: $s(u) = c + R u/\|u\|$;
  \item \textbf{Minkowski sum}: $s_{A \oplus B}(u) = s_A(u) + s_B(u)$.
\end{itemize}
The last identity is crucial: the support point of a Minkowski sum is the
sum of the support points of the summands, which is how GJK computes
support on the difference object $A \ominus B = A \oplus (-B)$ without ever
building it explicitly (Section~\ref{sec:convexint_minkowski_difference}).
Support functions also make the distance-to-origin queries of
Section~\ref{sec:convexint_distance} well defined for any convex body,
including those with smooth boundary, and they are the basis of the
Dobkin--Kirkpatrick hierarchy \cite{dobkin1983}.

\index{support function}\index{support point}

\section{Minkowski Difference}
\label{sec:convexint_minkowski_difference}

The \textbf{Minkowski difference} (or Minkowski sum of $A$ with $-B$) is
\[
  A \ominus B = A \oplus (-B) = \{ a - b : a \in A, b \in B\}.
\]
The Minkowski difference reformulates every intersection and distance
question about $A$ and $B$ as a question about the origin and
$A \ominus B$:

\begin{theorem}[Minkowski Difference Reductions]\label{thm:convexint_minkowski}
  For convex sets $A, B$:
  \begin{itemize}
    \item $A \cap B \neq \emptyset$ iff $\mathbf{0} \in A \ominus B$;
    \item $\operatorname{dist}(A, B) = \operatorname{dist}(\mathbf{0},
      A \ominus B)$;
    \item if $\mathbf{0} \notin A \ominus B$, the penetration depth of
      Section~\ref{sec:convexint_penetration_depth} is the distance from
      the origin to the boundary of $A \ominus B$.
  \end{itemize}
\end{theorem}

If $A$ and $B$ are convex, so is $A \ominus B$. The boundary of
$A \ominus B$ is made of translated copies of features of $A$ and $B$
(vertices, edges, faces), which is what GJK's simplex search exploits: the
closest point of $A \ominus B$ to the origin lies on a simplex whose
vertices are support points of $A \ominus B$, and that simplex is found by
iteratively adding the support point in the direction of the current
closest point and discarding redundant vertices. The
support-point construction of Section~\ref{sec:convexint_support_functions}
makes every evaluation of $A \ominus B$ cheap: $s_{A\ominus B}(u) =
s_A(u) - s_B(-u)$. This is the computational heart of GJK.

\index{Minkowski difference}

\section{Gilbert--Johnson--Keerthi Algorithm}
\label{sec:convexint_gjk}

\subsection{1. Overview}

The \textbf{Gilbert--Johnson--Keerthi (GJK) algorithm} \cite{gilbert1988}
computes the distance between two convex sets $A$ and $B$, or decides that
they intersect, using only support-function evaluations. Because it never
materializes the Minkowski difference, it runs in near-constant time for
simple primitives (boxes, spheres, convex hulls of a few points) and is the
standard narrow-phase test for convex objects in real-time physics engines
\cite{ericson2004,vandenbergen2004}.

\subsection{2. Definitions}

\begin{definition}[Simplex in Minkowski Space]\label{def:convexint_simplex}
  A \textbf{simplex} $\sigma \subseteq A \ominus B$ is a set of up to $d+1$
  support points $\{s(u_1), \dots, s(u_k)\}$, $k \le d+1$, generated by the
  algorithm. The origin's location relative to $\operatorname{conv}(\sigma)$
  determines whether $A$ and $B$ intersect.
\end{definition}

\subsection{3. Theory}

GJK iterates the following scheme. Start with an arbitrary support point of
$A \ominus B$ in some direction. At each step, let $v$ be the point of
$\operatorname{conv}(\sigma)$ closest to the origin. If $v = \mathbf{0}$,
then $\mathbf{0} \in A \ominus B$ and the objects intersect. Otherwise
compute $w = s_{A\ominus B}(-v)$, the support point in the direction
opposite the current closest point. If $\langle w, -v\rangle \le
\langle v, -v\rangle$ --- equivalently, if the origin is closer to
$\mathbf{0}$ than $w$ in direction $-v$ --- then no point of $A \ominus B$
is closer to the origin than $v$, and the distance is $\|v\|$. Otherwise
add $w$ to the simplex, remove any vertex not needed to keep the origin in
the affine hull of the simplex, and repeat.

\begin{theorem}[GJK Termination]\label{thm:convexint_gjk_term}
  The GJK iteration terminates: the distance of the current closest point
  to the origin strictly decreases at every step, and the stopping
  criterion is met when the support point in the test direction does not
  improve on the current closest point. In $\mathbb{R}^3$ the simplex has
  at most four vertices, so each iteration is $O(1)$.
\end{theorem}

The number of iterations is small in practice (typically a few dozen for
polyhedra); the worst case is bounded by the number of features but is not
polynomial in the input size, which is why GJK is used as an iterative
numerical algorithm rather than analyzed to an exact worst-case bound.

\subsection{4. Pseudocode}

\begin{algorithm}[htbp]
\caption{GJK Distance / Intersection Test}
\label{alg:convexint_gjk}
\begin{algorithmic}[1]
\Require Convex sets $A, B$ with support functions.
\Ensure Distance $\operatorname{dist}(A,B)$ and closest points, or ``intersect''.
\Function{GJK}{$A$, $B$}
  \State $d \gets (1,0,0)$; $S \gets \{s_{A\ominus B}(d)\}$
  \State $v \gets$ closest point of $\operatorname{conv}(S)$ to origin
  \While{true}
    \If{$\|v\| < \varepsilon$} \State \Return \textsc{Intersect} \EndIf
    \State $w \gets s_{A\ominus B}(-v)$
    \If{$\langle w, -v\rangle \le \langle v, v\rangle$}
      \State \Return $\|v\|$ (and the features supporting $v$)
    \EndIf
    \State $S \gets S \cup \{w\}$, reduced to the minimal simplex containing the origin in its affine hull
    \State $v \gets$ closest point of $\operatorname{conv}(S)$ to origin
  \EndWhile
\EndFunction
\end{algorithmic}
\end{algorithm}

\subsection{5. Parameter Selection}

\begin{itemize}
\item \textbf{Tolerance $\varepsilon$}: the termination threshold trades
  accuracy against iterations; for boolean queries a tolerance of $10^{-9}$
  in double precision is typical, and the convergence is quadratic once the
  simplex is right (the closest point approaches the origin).
\item \textbf{Initial direction}: any nonzero direction works; using the
  relative displacement of the centers of mass speeds up convergence.
\item \textbf{Caching}: caching the last support feature (warm starting)
  reduces iterations dramatically for nearby queries between the same
  objects across time steps.
\end{itemize}

\subsection{6. Complexity}

\begin{itemize}
\item \textbf{Per iteration}: $O(1)$ support evaluations (each $O(1)$ for
  simple primitives, $O(\log n)$ for a polygon, $O(n)$ for a general
  polyhedron), plus a constant-size closest-point computation on the
  simplex.
\item \textbf{Iterations}: a small constant in practice; the theoretical
  worst case is larger but does not arise on typical inputs.
\end{itemize}

\subsection{7. Usages}

\begin{itemize}
\item \textbf{Real-time physics}: the narrow phase of collision detection
  for convex bodies in physics engines such as PhysX, Havok, and Bullet
  \cite{ericson2004,vandenbergen2004}.
\item \textbf{Robotics}: distance queries in configuration spaces
  (Chapter~\ref{ch:motion_planning}).
\item \textbf{Haptics and surgical simulation}: fast contact queries for
  force feedback.
\end{itemize}

\subsection{8. Testing Methods and Edge Cases}

\begin{itemize}
\item \textbf{Touching objects}: verify the zero-distance threshold
  distinguishes touching from overlapping consistently.
\item \textbf{Deep penetration}: GJK with an origin strictly inside
  $A \ominus B$ returns ``intersect'' but no depth; the penetration depth
  of Section~\ref{sec:convexint_penetration_depth} requires EPA.
\item \textbf{Support correctness}: the support oracle must be exact on
  flat faces (all vertices of a face are support points); the choice of
  representative affects the closest-point computation.
\end{itemize}

\subsection{9. References}

GJK was introduced in \cite{gilbert1988}; the standard engineering
treatments with implementation details are \cite{ericson2004,vandenbergen2004}.

\section{Penetration Depth}
\label{sec:convexint_penetration_depth}

\subsection{1. Overview}

When two convex objects overlap, the \textbf{penetration depth} measures how
deeply: the minimum translation that separates them. For convex $A, B$, the
penetration depth is the distance from the origin to the boundary of the
Minkowski difference $A \ominus B$ (Theorem~\ref{thm:convexint_minkowski}),
and the separating translation is the closest boundary point. The
\textbf{expanding polytope algorithm (EPA)}\index{expanding polytope algorithm}
computes this distance by iteratively expanding a polytope inside
$A \ominus B$ toward the boundary point closest to the origin
\cite{vandenbergen2004}.

\subsection{2. Definitions}

\begin{definition}[Penetration Depth]\label{def:convexint_penetration_depth}
  The \textbf{penetration depth} of overlapping convex sets $A, B$ is
  $\min\{ \|t\| : (A + t) \cap B = \emptyset \}$, the length of the
  smallest translation that separates them. The \textbf{penetration
  vector} is the translation achieving the minimum.
\end{definition}

\subsection{3. Theory}

By Theorem~\ref{thm:convexint_minkowski}, the penetration depth equals
$\operatorname{dist}(\mathbf{0}, \partial(A \ominus B))$. Since the
boundary of $A \ominus B$ is composed of facets whose support directions are
the feature normals, the closest boundary point is found by expanding a
simplex (from GJK) into a polytope: maintain the polytope's facets, find
the facet closest to the origin, push the polytope outward at the support
point in the facet's normal direction, and repeat until the closest facet
converges. EPA is correct because the polytope is always contained in
$A \ominus B$ and expands toward the closest boundary point.

\subsection{4. Pseudocode}

\begin{algorithm}[htbp]
\caption{Expanding Polytope Algorithm (Penetration Depth)}
\label{alg:convexint_epa}
\begin{algorithmic}[1]
\Require Simplex $S$ from GJK with origin in $\operatorname{conv}(S) \subseteq A \ominus B$.
\Ensure Penetration vector of $A, B$.
\Function{EPA}{$A$, $B$, $S$}
  \State build a polytope $P$ (triangulation of the boundary of $S$)
  \While{true}
    \State $f \gets$ facet of $P$ closest to the origin; $n \gets$ normal of $f$
    \State $w \gets s_{A\ominus B}(n)$
    \If{$\langle w, n\rangle - \langle \text{point of } f, n\rangle < \varepsilon$}
      \State \Return $\langle \text{point of } f, n\rangle \cdot n$  \Comment{penetration vector}
    \EndIf
    \State add $w$ to $P$; retriangulate the visible facets
  \EndWhile
\EndFunction
\end{algorithmic}
\end{algorithm}

\subsection{5. Parameter Selection}

\begin{itemize}
\item \textbf{Initial simplex}: EPA starts from the GJK simplex containing
  the origin; the polytope quality determines convergence speed.
\item \textbf{Convergence threshold}: the $\varepsilon$ on the facet
  support controls accuracy versus iterations; for game physics a value
  near $10^{-4}$ of the object scale is common.
\end{itemize}

\subsection{6. Complexity}

\begin{itemize}
\item \textbf{Time}: $O(1)$ per iteration plus $O(1)$ support evaluations;
  the polytope grows linearly in the number of iterations, and the
  iteration count is a small constant in practice.
\item \textbf{Space}: $O(k)$ for the polytope with $k$ facets.
\end{itemize}

\subsection{7. Usages}

\begin{itemize}
\item \textbf{Collision response}: computing the separating impulse
  magnitude in physics engines.
\item \textbf{Contact manifolds}: the penetration vector defines the
  contact normal for solver-based dynamics.
\end{itemize}

\subsection{8. Testing Methods and Edge Cases}

\begin{itemize}
\item \textbf{Deep overlaps}: verify convergence when the origin is far
  inside $A \ominus B$.
\item \textbf{Flat faces}: a facet parallel to a face of the difference
  gives a unique closest point; test boxes and degenerate (flat) convex
  bodies.
\item \textbf{Consistency}: for small overlaps, the penetration vector must
  agree with the vector computed by a reference translation test.
\end{itemize}

\subsection{9. References}

EPA is described in the game-physics literature
\cite{vandenbergen2004,ericson2004}; its Minkowski-difference foundation is
the classical theory of convex sets (Section~\ref{sec:math_separating_surfaces}).

\section{Collision Detection}
\label{sec:convexint_collision}

Collision detection in an interactive or physical system is organized in
two phases. The \textbf{broad phase} prunes pairs of objects that cannot
possibly touch using cheap bounding volumes: each object is wrapped in an
axis-aligned bounding box (AABB), and pairs whose boxes do not overlap are
discarded. The canonical broad-phase algorithms are:

\begin{description}
  \item[Sweep and prune] sort the boxes by one coordinate and slide a
    window, testing only pairs within the window; for moderately sized
    scenes this is near-linear because the window is small
    \cite{ericson2004}.
  \item[Boundary volume hierarchies] each object stores a tree of
    bounding volumes (boxes or spheres); overlapping subtrees are descended
    recursively, giving output-sensitive behavior for complex objects
    (Chapter~\ref{ch:spatial_structures}).
  \item[Uniform and spatial hashing] objects are bucketed into cells of a
    grid; only pairs sharing a cell are tested, which suits uniformly
    distributed scenes.
\end{description}

The \textbf{narrow phase} then tests the surviving pairs exactly. For
convex pairs, the narrow phase uses SAT (Section~\ref{sec:convexint_sat})
for boxes and GJK (Section~\ref{sec:convexint_gjk}) for general convex
bodies, followed by EPA when penetration is needed
(Section~\ref{sec:convexint_penetration_depth}). The overall cost is
dominated by the broad phase, so the engineering advice is to make the
bounding volumes tight, keep the broad-phase prune aggressive, and spend
the narrow-phase budget only on pairs that truly overlap. The spatial data
structures of Chapter~\ref{ch:spatial_structures} and the sweep techniques
of Chapter~\ref{ch:segment_intersection} provide the primitives.

\index{collision detection}\index{broad phase}\index{narrow phase}

\section{Continuous Collision Detection}
\label{sec:convexint_continuous}

\subsection{1. Overview}

\textbf{Continuous collision detection} (CCD)\index{continuous collision
detection} asks not only whether two objects overlap at the current time,
but whether they overlap at \emph{any} time in an interval --- and, if so,
the first time of contact (the \textbf{time of impact}, TOI). Discrete
tests at the endpoints of the interval can miss a collision that occurs
between the sampled instants (tunneling), so CCD computes the swept
behavior.

\subsection{2. Definitions}

\begin{definition}[Time of Impact]\label{def:convexint_toi}
  For bodies moving on prescribed trajectories during $[0,1]$, the
  \textbf{time of impact} is the smallest $t$ at which two bodies touch.
\end{definition}

\subsection{3. Theory}

The standard CCD method for convex objects is \textbf{conservative
advancement}: at the current time, compute the distance $\delta$ between
the two bodies (with GJK, Section~\ref{sec:convexint_gjk}) and bound the
relative speed $v_{\max}$ over the interval; then no collision can occur
before the objects close the distance, so the time is advanced by
$\delta / v_{\max}$ and the process repeats until the step is below a
tolerance. Conservative advancement is correct as long as the speed bound
is valid, which makes it the default for convex bodies. An alternative for
polyhedra is the exact swept-volume/feature-interval method: subdivide time
at the events where the separating feature changes, and test each interval
algebraically \cite{ericson2004}.

\subsection{4. Pseudocode}

\begin{algorithm}[htbp]
\caption{Conservative Advancement (Time of Impact)}
\label{alg:convexint_ccd}
\begin{algorithmic}[1]
\Require Convex bodies $A(t), B(t)$, trajectories over $[0,1]$, tolerance $\varepsilon$.
\Ensure The time of impact $t$ (or no collision).
\Function{ConservativeAdvancement}{$A$, $B$}
  \State $t \gets 0$
  \While{$t \le 1$}
    \State $\delta \gets \textsc{GJK}(A(t), B(t))$
    \If{$\delta \le \varepsilon$} \State \Return $t$ \EndIf
    \State $v_{\max} \gets$ bound on relative speed of closest features over $[t,1]$
    \State $t \gets \min(1, t + \delta / v_{\max})$
  \EndWhile
  \State \Return \textsc{NoCollision}
\EndFunction
\end{algorithmic}
\end{algorithm}

\subsection{5. Parameter Selection}

\begin{itemize}
\item \textbf{Speed bound}: a global Lipschitz bound on the relative motion
  makes the algorithm conservative (never misses a collision) but may step
  slowly; a local bound on the closest features is tighter.
\item \textbf{Tolerance}: $\varepsilon$ controls the accuracy of the
  reported TOI; the cost grows logarithmically with $1/\varepsilon$.
\end{itemize}

\subsection{6. Complexity}

\begin{itemize}
\item \textbf{Time}: $O(I \cdot T)$ where $T$ is the GJK cost and $I$ the
  number of advancement steps, which is $O(\log(\text{scale}/\varepsilon))$
  for linearly separated trajectories.
\end{itemize}

\subsection{7. Usages}

\begin{itemize}
\item \textbf{Games and physics}: preventing tunneling of fast bullets and
  thin objects.
\item \textbf{Robotics}: verifying that a moving manipulator path is
  collision-free (Chapter~\ref{ch:motion_planning}).
\end{itemize}

\subsection{8. Testing Methods and Edge Cases}

\begin{itemize}
\item \textbf{Tunneling regression}: a fast small object passing through a
  thin wall must be reported, not missed.
\item \textbf{Graze and rest contact}: contact that lasts exactly at an
  instant (tangent) must be returned consistently.
\item \textbf{Consistency}: compare TOI against fine-grained discrete
  sampling on random motions.
\end{itemize}

\subsection{9. References}

CCD and conservative advancement are treated in
\cite{ericson2004,vandenbergen2004}.

\section{Applications in Robotics and Simulation}
\label{sec:convexint_applications}

The intersection and distance algorithms of this chapter power the
collision layer of three industries:

\begin{description}
  \item[Game and simulation physics] Real-time physics engines ---
    NVIDIA PhysX, Havok, and Bullet --- run SAT, GJK, and EPA millions of
    times per second to maintain non-penetrating rigid-body simulations; the
    algorithms of this chapter (broad-phase bounding hierarchies, SAT for
    boxes, GJK for convex hulls, EPA for penetration, and conservative
    advancement for CCD) are the exact components they ship
    \cite{ericson2004,vandenbergen2004}.
  \item[Robotics] Motion planners compute distances between the robot's
    links and its environment using GJK-based distance queries over convex
    decompositions of non-convex links, and the configuration-space
    techniques of Chapter~\ref{ch:motion_planning} rely on these queries
    for feasibility tests. Collision-free programming for industrial arms
    is a direct consumer.
  \item[CAD and CAM] Clearance analysis, toolpath validation, and assembly
    feasibility checks are distance and penetration queries between convex
    primitives and convex hulls of tooling geometry; the closest-feature
    results are reused for tolerancing and clearance reporting.
\end{description}

In each setting the convex intersection chapter's contribution is the
same: a small set of primitives (support oracles, projection tests, and
simplex geometry) from which the entire collision pipeline is assembled,
with the exactness and robustness requirements of
Sections~\ref{sec:intro_robustness}--\ref{sec:intro_exact_computation}.

\index{physics engine}\index{collision}

\section{Exercises}
\label{sec:convexint_exercises}

\begin{enumerate}
  \item Prove the separating axis theorem for two convex polygons from the
    hyperplane-separation theorem of Section~\ref{sec:math_separating_surfaces}.
  \item Show that the projection interval of a convex polygon onto an axis
    is determined by its support function values at the axis and its
    negation.
  \item Verify the Minkowski-difference reductions of
    Theorem~\ref{thm:convexint_minkowski} for two boxes, one contained in
    the other.
  \item Implement GJK for two polygons in the plane and verify it against a
    brute-force distance computation on random inputs.
  \item Prove that in $\mathbb{R}^2$ the GJK simplex has at most three
    vertices, and that the stopping test $\langle w, -v\rangle \le \langle
    v, v\rangle$ is necessary and sufficient.
  \item Explain why GJK alone cannot report penetration depth, and
    construct an input where the origin lies strictly inside
    $A \ominus B$.
  \item Adapt the separating axis test to axis-aligned boxes in
    $\mathbb{R}^3$ and show it requires only three interval comparisons per
    candidate axis.
  \item Simulate a tunneling failure: sample two trajectories at their
    endpoints and show a missed collision that conservative advancement
    (Section~\ref{sec:convexint_continuous}) detects.
\end{enumerate}
